\documentclass[11pt]{article}

\usepackage[T1]{fontenc}
\usepackage{amsmath,amssymb,amsthm,mathtools}
\usepackage[margin=1in]{geometry}
\usepackage[colorlinks=true,linkcolor=blue,citecolor=blue,urlcolor=blue]{hyperref}

\newtheorem{theorem}{Theorem}
\newtheorem{lemma}[theorem]{Lemma}
\newtheorem{proposition}[theorem]{Proposition}
\theoremstyle{definition}
\newtheorem{definition}[theorem]{Definition}
\theoremstyle{remark}
\newtheorem{remark}[theorem]{Remark}

\newcommand{\R}{\mathbb R}
\newcommand{\Q}{\mathbb Q}
\newcommand{\A}{\mathcal A}
\newcommand{\V}{\mathcal V}
\newcommand{\dd}{\mathrm d}
\newcommand{\grad}{\operatorname{grad}}
\newcommand{\inner}[2]{\left\langle #1,#2\right\rangle_Q}
\newcommand{\normQ}[1]{\left\lVert #1\right\rVert_Q}

\title{An Explicit Affine-Chart Proof of\\
Veronese Penalty Coercivity\\
\large Companion to the Lean formalization}
\author{}
\date{\today}

\begin{document}
\maketitle

\begin{abstract}
We give a self-contained proof of the local inequality
\[
  \normQ{\grad_{S_Q}W(y)}^2\geq c\,W(y)
\]
for the complete sum-of-squares penalty of the degree-\(k\) Veronese
cone.  The usual proof invokes the Morse--Bott lemma and a tubular
neighborhood.  Here we instead use an explicit affine reconstruction of
the cone and a normalized bad-sequence argument.  This is both elementary
and close to the accompanying Lean development.  The proof also shows
that the tube radius and coercivity constant may be chosen to be the
same reciprocal integer.
\end{abstract}

\section{The penalty and its spherical gradient}

Fix integers \(n\geq 0\) and \(k\geq 1\).  Let
\[
  \A_k=\left\{\alpha\in\mathbb N^n:
      |\alpha|:=\sum_{\ell=1}^n\alpha_\ell=k\right\},
  \qquad E=\R^{\A_k}.
\]
The degree-\(k\) Veronese map is
\[
  \nu_k:\R^n\longrightarrow E,\qquad
  \nu_k(x)_\alpha=x^\alpha.
\]
Put
\[
  Q(y)=\sum_{\alpha\in\A_k}
      {k\choose\alpha}y_\alpha^2,
  \qquad
  \inner uv=\sum_{\alpha\in\A_k}
      {k\choose\alpha}u_\alpha v_\alpha.
\]
Thus \(Q(y)=\normQ y^2\), and \(Q\) is positive definite.

A \emph{balanced quadratic relation} is a quadruple
\(r=(\alpha,\beta,\gamma,\delta)\in\A_k^4\) with
\(\alpha+\beta=\gamma+\delta\).  Write
\[
  q_r(y)=y_\alpha y_\beta-y_\gamma y_\delta.
\]
We use every balanced quadruple; redundancy is harmless and makes the
family completely explicit.  Its common zero set is the real affine
Veronese cone
\[
  \V=\{y:q_r(y)=0\text{ for every }r\}.
\]
The complete constraint energy is
\[
  W(y)=\sum_r q_r(y)^2.
\]
In particular,
\[
  W(y)\geq0,\qquad W(y)=0\Longleftrightarrow y\in\V.
\]

Let \(S_Q=\{y:Q(y)=1\}\).  Denote the ambient gradient for the
\(Q\)-metric by \(\grad_QW\).  Since \(W\) is homogeneous of degree
four, Euler's identity gives
\[
  \grad_{S_Q}W(y)=\grad_QW(y)-4W(y)y
  \quad (y\in S_Q).
\]
Set
\[
  D(y)=\normQ{\grad_{S_Q}W(y)}^2.
\]

\begin{theorem}[Veronese penalty coercivity]\label{thm:coercivity}
There is an integer \(m\geq0\) such that, with
\[
  \varepsilon_m=\frac1{m+1},
\]
one has
\[
  Q(y)=1,\qquad W(y)\leq\varepsilon_m
  \quad\Longrightarrow\quad
  \varepsilon_m W(y)\leq D(y).
\]
Consequently the local Morse--Bott inequality holds with positive
rational constants \(c=\eta=\varepsilon_m\).
\end{theorem}

\section{An explicit affine reconstruction}

Fix a coordinate \(i\in\{1,\dots,n\}\).  Let
\[
  b=ke_i,\qquad
  b_j=(k-1)e_i+e_j.
\]
When \(y_b\neq0\), define affine source coordinates
\[
  x_j(y)=\frac{y_{b_j}}{y_b}
\]
and reconstruct a cone point by
\[
  C_i(y)=y_b\,\nu_k(x(y)).
  \tag{1}\label{eq:chart}
\]

\begin{lemma}[Chart identities]\label{lem:chart}
Whenever \(y_b\neq0\), the reconstruction \(C_i(y)\) has the following
properties.
\begin{enumerate}
\item \(C_i(y)\in\V\).
\item \(C_i(y)_{b_j}=y_{b_j}\) for every \(j\), including \(j=i\),
      for which \(b_i=b\).
\item If \(y\in\V\), then \(C_i(y)=y\).
\item The map \(C_i\) is continuous at every point where \(y_b\neq0\).
\end{enumerate}
\end{lemma}

\begin{proof}
The first assertion follows because \(C_i(y)\) is a scalar multiple of
a Veronese point and every \(q_r\) is homogeneous quadratic.  For the
second, \(x_i(y)=1\), and hence
\[
  \nu_k(x(y))_{b_j}=x_i(y)^{k-1}x_j(y)=x_j(y).
\]
Multiplication by \(y_b\) recovers \(y_{b_j}\).

For the third assertion, repeatedly use the balanced relation
\[
  \alpha+ke_i
   =(\alpha+e_i-e_j)+((k-1)e_i+e_j)
  \tag{2}\label{eq:transfer}
\]
whenever \(\alpha_j>0\) and \(j\neq i\).  It determines \(y_\alpha\)
from a coordinate with one more unit in its \(i\)-th exponent.  Induction
on \(k-\alpha_i\) eventually reaches the pure coordinate \(ke_i\), and
gives
\[
  y_\alpha=y_b\,x(y)^\alpha=C_i(y)_\alpha.
\]
Finally, every coordinate in \eqref{eq:chart} is a finite product of
powers of quotients with common denominator \(y_b\), so it is continuous
where that denominator is nonzero.
\end{proof}

Define the chart error
\[
  e_i(y)=y-C_i(y).
\]
By Lemma~\ref{lem:chart},
\[
  e_i(y)_{b_j}=0\qquad\text{for every }j.
  \tag{3}\label{eq:gauge}
\]
Moreover, \(e_i(y)=0\) implies \(y=C_i(y)\in\V\), so
\[
  W(y)>0\quad\Longrightarrow\quad e_i(y)\neq0.
  \tag{4}\label{eq:error-nonzero}
\]

\section{The chart-gauged differential kernel}

The differential of \(q_r\) at \(p\), in direction \(h\), is denoted
\(\dd q_r(p)h\).

\begin{lemma}[Trivial chart-gauged kernel]\label{lem:kernel}
Let \(p,h\in E\), suppose \(p_b\neq0\), and assume
\[
  h_{b_j}=0\quad\text{for every }j,
  \qquad
  \dd q_r(p)h=0\quad\text{for every balanced }r.
\]
Then \(h=0\).
\end{lemma}

\begin{proof}
We prove \(h_\alpha=0\) by strong induction on
\(d(\alpha)=k-\alpha_i\).  If \(d(\alpha)=0\), then
\(\alpha=ke_i=b=b_i\), and the result is one of the chart-coordinate
assumptions.

Suppose \(d(\alpha)>0\).  Some \(j\neq i\) has \(\alpha_j>0\).  Put
\[
  \alpha'=\alpha+e_i-e_j.
\]
The defect of \(\alpha'\) is one smaller, so \(h_{\alpha'}=0\) by
induction.  Apply the differential hypothesis to the transfer relation
\eqref{eq:transfer}.  It reads
\[
  h_\alpha p_b+p_\alpha h_b
   -h_{\alpha'}p_{b_j}-p_{\alpha'}h_{b_j}=0.
\]
The last three terms vanish: \(h_b=h_{b_j}=0\) by the chart gauge, and
\(h_{\alpha'}=0\) by induction.  Thus \(h_\alpha p_b=0\).  Since
\(p_b\neq0\), we obtain \(h_\alpha=0\).
\end{proof}

\begin{remark}
No smooth-manifold, tangent-space, or scheme-theoretic statement is
hidden in Lemma~\ref{lem:kernel}.  The point \(p\) is not even required
to lie in \(\V\).  The lemma is a finite coordinate induction using the
specific balanced quadrics \eqref{eq:transfer}.
\end{remark}

\section{Proof of coercivity}

\begin{proof}[Proof of Theorem~\ref{thm:coercivity}]
Assume the conclusion is false.  For every \(m\geq0\), choose \(y_m\)
such that
\[
  Q(y_m)=1,\qquad
  W(y_m)\leq\varepsilon_m,\qquad
  D(y_m)<\varepsilon_m W(y_m).
  \tag{5}\label{eq:bad}
\]
Since \(D\geq0\), the last strict inequality implies \(W(y_m)>0\).
Also \(\varepsilon_m\to0\), hence \(W(y_m)\to0\).

\paragraph{First compactness step.}
The ellipsoid \(S_Q\) is compact.  After passing to a subsequence,
\[
  y_m\longrightarrow p\in S_Q.
\]
Continuity gives \(W(p)=0\), so \(p\in\V\).  Since \(Q(p)=1\), the point
\(p\) is nonzero.  A nonzero point of the real Veronese cone has a
nonzero pure coordinate; choose \(i\) so that \(p_b\neq0\), with
\(b=ke_i\).  After discarding finitely many terms, \(y_{m,b}\neq0\) for
every \(m\).

\paragraph{Normalize the chart error.}
Write
\[
  C_m=C_i(y_m),\qquad e_m=y_m-C_m,\qquad
  \rho_m=\sqrt{Q(e_m)},\qquad v_m=\rho_m^{-1}e_m.
  \tag{6}\label{eq:normalize}
\]
Because \(W(y_m)>0\), equation \eqref{eq:error-nonzero} shows
\(e_m\neq0\), so \(\rho_m>0\).  The normalization gives
\[
  Q(v_m)=1,\qquad (v_m)_{b_j}=0\quad\text{for every }j.
  \tag{7}\label{eq:normalized-properties}
\]
The chart is continuous at \(p\) and fixes \(p\).  Therefore
\[
  C_m\to p,\qquad e_m\to0,\qquad \rho_m\to0.
\]
Compactness of \(S_Q\) gives a second subsequence for which
\[
  v_m\longrightarrow v,\qquad
  Q(v)=1,\qquad v_{b_j}=0\quad\text{for every }j.
  \tag{8}\label{eq:v-limit}
\]

\paragraph{Normalize the quadratic relations.}
For each balanced relation \(r\), put
\[
  A_{m,r}=\frac{q_r(y_m)}{\rho_m},
  \qquad
  a_r=\dd q_r(p)v.
\]
The exact quadratic Taylor formula is
\[
  q_r(C_m+e_m)
   =q_r(C_m)+\dd q_r(C_m)e_m+q_r(e_m).
\]
Here \(q_r(C_m)=0\) and \(e_m=\rho_m v_m\).  By linearity of the
differential in its direction and quadratic homogeneity of \(q_r\),
\[
  A_{m,r}
   =\dd q_r(C_m)v_m+\rho_m q_r(v_m)
   \longrightarrow \dd q_r(p)v=a_r.
  \tag{9}\label{eq:A-limit}
\]
Let
\[
  S_m=\sum_r A_{m,r}^2,\qquad S=\sum_r a_r^2.
\]
The relation family is finite, so \(S_m\to S\).  Directly from the
definition,
\[
  W(y_m)=\rho_m^2S_m.
  \tag{10}\label{eq:W-factor}
\]

\paragraph{The normalized spherical gradient tends to zero.}
Let
\[
  s_m=\grad_{S_Q}W(y_m),\qquad t_m=\rho_m^{-1}s_m.
\]
Because \(Q(s_m)=D(y_m)\), equations
\eqref{eq:bad} and \eqref{eq:W-factor} imply
\[
  Q(t_m)=\frac{D(y_m)}{\rho_m^2}
  <\varepsilon_m S_m\longrightarrow0.
  \tag{11}\label{eq:t-zero}
\]
Weighted Cauchy--Schwarz and \(Q(v_m)=1\) now give
\[
  |\inner{t_m}{v_m}|^2
  \leq Q(t_m)Q(v_m)=Q(t_m),
\]
so
\[
  \inner{t_m}{v_m}\longrightarrow0.
  \tag{12}\label{eq:t-pairing}
\]

\paragraph{Compute the same limit using the ambient gradient.}
The ambient and spherical gradients satisfy
\[
  \grad_QW(y_m)=s_m+4W(y_m)y_m.
\]
Divide its pairing with \(v_m\) by \(\rho_m\).  Using
\eqref{eq:W-factor},
\[
 \begin{aligned}
  \frac{\inner{\grad_QW(y_m)}{v_m}}{\rho_m}
   &=
     \inner{t_m}{v_m}
     +4\rho_m S_m\inner{y_m}{v_m}\\
   &\longrightarrow0.
 \end{aligned}
 \tag{13}\label{eq:ambient-zero}
\]
Indeed, the first term tends to zero by \eqref{eq:t-pairing}; in the
second, \(\rho_m\to0\), while \(S_m\) and
\(\inner{y_m}{v_m}\) converge and are therefore bounded.

On the other hand, differentiating the sum of squares gives the exact
identity
\[
  \inner{\grad_QW(y_m)}{v_m}
   =2\sum_r q_r(y_m)\,\dd q_r(y_m)v_m.
\]
After division by \(\rho_m\), equations \eqref{eq:A-limit} and
\(y_m\to p,\ v_m\to v\) show that the left side of
\eqref{eq:ambient-zero} also converges to
\[
  2\sum_r a_r^2=2S.
\]
Consequently \(S=0\), and hence
\[
  \dd q_r(p)v=a_r=0\qquad\text{for every }r.
\]

Finally, \(p_b\neq0\), all chart coordinates of \(v\) vanish by
\eqref{eq:v-limit}, and all relation differentials annihilate \(v\).
Lemma~\ref{lem:kernel} gives \(v=0\), contradicting \(Q(v)=1\).
This contradiction proves the theorem.
\end{proof}

\section{Rational effectivity and the formalization boundary}

Theorem~\ref{thm:coercivity} already produces positive rational
constants: \(c=\eta=1/(m+1)\).  If one needs to find a valid pair from
the polynomial data, enumerate positive rational pairs \((c,\eta)\) and
decide the first-order sentence
\[
  \forall y\in E,\quad
  Q(y)=1\ \wedge\ 0\leq W(y)\leq\eta
  \ \Longrightarrow\ cW(y)\leq D(y).
  \tag{14}\label{eq:qe}
\]
All expressions in \eqref{eq:qe} are rational polynomials, so the
sentence is decidable by quantifier elimination for real closed fields;
see, for example, \cite[Chapters 1 and 2]{BCR}.  Existence from
Theorem~\ref{thm:coercivity} guarantees that the enumeration terminates.

The Lean development formalizes the full analytic existence argument,
including compactness, both subsequence extractions, every exact
quadratic identity, and the final kernel contradiction.  The generic
real-closed-field quantifier-elimination implementation used to turn
\eqref{eq:qe} into an executable search is a separate computability
interface; it is not silently assumed inside the coercivity theorem.

\section{Guide to the Lean modules}

The mathematical steps above correspond to the following reusable
modules and declarations.
\begin{itemize}
\item \texttt{HGAS.Veronese.Conormal} packages the differential of the
      Veronese map and the pure/adjacent coordinate identities.
\item \texttt{HGAS.Veronese.AffineConeChart} defines
      \texttt{affineConeChart} and \texttt{affineConeError}, proves the
      chart identities and continuity, proves the exact quadratic
      expansion, establishes weighted Cauchy--Schwarz, and proves
      \texttt{eq\_zero\_of\_toricRelationDifferential\_eq\_zero\_of\_edge\_eq\_zero}.
\item \texttt{HGAS.Veronese.MorseBottCoercivity} proves the normalized
      bad-sequence contradiction and exposes the reciprocal-integer and
      positive-constant coercivity theorems.
\item \texttt{HGAS.Veronese.CubicLiftToricFactorization} supplies the
      independent \(O(W)\) tangency estimate for the cubic lift by exact
      ideal factorization.
\item \texttt{HGAS.Veronese.TransverseStabilization} combines these two
      estimates with the exact stabilization identities.
\end{itemize}

\begin{thebibliography}{9}

\bibitem{BCR}
J.~Bochnak, M.~Coste, and M.-F.~Roy,
\emph{Real Algebraic Geometry},
Ergebnisse der Mathematik und ihrer Grenzgebiete, vol.~36,
Springer, 1998.

\end{thebibliography}

\end{document}
